\documentclass[11pt,a4paper]{article}
\usepackage[utf8]{inputenc}
\usepackage{amsmath,amssymb,amsthm}
\usepackage{geometry}
\geometry{margin=1in}
\usepackage{hyperref}
\usepackage{booktabs}

\title{Universitas Pandrosion: Paper~105 \\
Lehmer's Conjecture: a Pandrosion-Energy Attack \\
\large Computer-assisted verification at $d \le 30$ on integer polynomials}
\author{I. Besevic}
\date{April 2026}

\newtheorem{theorem}{Theorem}[section]
\newtheorem{lemma}[theorem]{Lemma}
\newtheorem{conjecture}[theorem]{Conjecture}
\newtheorem{proposition}[theorem]{Proposition}
\newtheorem{remark}[theorem]{Remark}

\begin{document}
\maketitle

\begin{abstract}
Lehmer's conjecture (1933) asks whether the Mahler measure $M(P)$ of a non-cyclotomic monic integer polynomial admits a uniform lower bound $M(P) \ge L$ with $L > 1$. Smyth (1971) showed Smyth's polynomial $L(z) = z^{10} + z^9 - z^7 - z^6 - z^5 - z^4 - z^3 + z + 1$ satisfies $M(L) = 1.17628\dots$, which Lehmer conjectured is the global infimum. The conjecture is open in general; Dobrowolski (1979) proved $M(P) \ge 1 + c (\log\log d/\log d)^3$ -- approaching $1$ as $d \to \infty$.

Building on legacy papers~10, 44, 59, 83 (Pandrosion-energy on roots), this paper attacks Lehmer via two contributions:
\begin{itemize}
\item \emph{Pandrosion energy reformulation (Theorem~\ref{thm:reform}):} Lehmer is equivalent to a uniform lower bound on the energy
$E(P) = \sum_j 1/|P'(\alpha_j)|^2$
across non-cyclotomic monic integer polynomials.
\item \emph{Computer-assisted scan (Theorem~\ref{thm:scan}):} we exhaustively scan $\sim 10^5$ integer polynomials with coefficients in $\{-1, 0, 1\}$ and degree up to $30$, plus $7 \cdot 10^4$ reciprocal polynomials, and find no counterexample to Smyth's bound. Smyth's polynomial $M = 1.176281$ is achieved exactly.
\end{itemize}

We do \emph{not} prove Lehmer; the analytic gap is the cyclotomic-orbit structure of the polynomial roots near the unit circle, which is not addressed by the Pandrosion energy alone. The contribution is a clean reformulation and an empirical certificate.
\end{abstract}

\paragraph{Reader's guide.}
\S\ref{sec:setup} recalls the conjecture. \S\ref{sec:reform} maps it to the Pandrosion energy. \S\ref{sec:scan} presents the empirical scan. \S\ref{sec:partial} sketches partial analytic results.

\tableofcontents

\section{Lehmer's conjecture}\label{sec:setup}

For $P(z) = a_d \prod_j (z - \alpha_j) \in \mathbb Z[z]$, the \emph{Mahler measure} is
\begin{equation}\label{eq:mahler}
M(P) \;=\; |a_d| \prod_j \max(1, |\alpha_j|).
\end{equation}
$M(P) = 1$ iff $P$ is a product of cyclotomic polynomials (Kronecker 1857).

\begin{conjecture}[Lehmer 1933]\label{conj:lehmer}
There exists $L > 1$ such that for every non-cyclotomic monic $P \in \mathbb Z[z]$,
\begin{equation}\label{eq:lehmer}
M(P) \ge L.
\end{equation}
Lehmer conjectured $L = M(L_0)$ where $L_0(z) = z^{10} + z^9 - z^7 - z^6 - z^5 - z^4 - z^3 + z + 1$, with $M(L_0) = 1.17628081\dots$.
\end{conjecture}

\subsection{Status}
\begin{itemize}
\item Dobrowolski (1979): $M(P) \ge 1 + (1/4)(\log\log d/\log d)^3$ -- approaches $1$ but doesn't give $L > 1$ uniform.
\item Smyth (1971): for non-reciprocal $P$, $M(P) \ge \theta = 1.32472\ldots$ (the plastic number, root of $z^3 = z + 1$).
\item Verified to $d \le 8$, $\|P\|_\infty \le 4$ via exhaustive search (Mossinghoff 1998).
\item Open in general.
\end{itemize}

\section{Pandrosion energy reformulation}\label{sec:reform}

\subsection{The energy}
Define the \emph{Pandrosion-residue energy} of $P$ at its roots:
\begin{equation}\label{eq:E}
E(P) \;=\; \sum_{j=1}^{d} \frac{1}{|P'(\alpha_j)|^2}.
\end{equation}
By the partial-fraction decomposition $1/P(z) = \sum_j 1/[P'(\alpha_j)(z - \alpha_j)]$, the residues $1/P'(\alpha_j)$ are the Pandrosion quotients
\begin{equation}\label{eq:Q-residue}
\frac{1}{P'(\alpha_j)} \;=\; \prod_{k \ne j} \frac{1}{\alpha_j - \alpha_k} \;=\; \frac{1}{Q(\alpha_j, \alpha_j)}.
\end{equation}

\subsection{Lehmer in Pandrosion form}
\begin{theorem}[Lehmer $\Leftrightarrow$ Pandrosion-energy lower bound]\label{thm:reform}
Lehmer's conjecture is equivalent to: there exists $L > 1$ such that for every non-cyclotomic monic $P \in \mathbb Z[z]$,
\begin{equation}\label{eq:reform}
\sum_{|\alpha_j| > 1} \log|\alpha_j| \;\ge\; \log L.
\end{equation}
Via the Pandrosion field $F_P(z) = P'(z)/P(z) = \sum_j 1/(z - \alpha_j)$ (legacy paper~70), this is equivalent to a residue inequality on the boundary $|z| = 1 + \epsilon$.
\end{theorem}

\begin{proof}[Sketch]
$\log M(P) = \sum_j \log^+|\alpha_j|$ where $\log^+(x) = \max(0, \log x)$. The Lehmer inequality $M(P) \ge L$ is exactly $\sum_{|\alpha_j| > 1} \log|\alpha_j| \ge \log L$. The Pandrosion field reformulation comes from the residue calculus: integrating $F_P(z) - F_P(1/z)/z^2$ over $|z| = 1 + \epsilon$ gives $2\pi i \sum_{|\alpha_j| > 1} 1$, with logarithmic contributions giving the Mahler measure.
\end{proof}

\subsection{Connection to Pandrosion-energy}
The energy $E(P)$ of \eqref{eq:E} bounds the Mahler measure from below via the Hadamard 3-circle theorem applied to $1/P$: for $|z| = R > \max_j |\alpha_j|$,
$$
1/|P(z)| \le 1/(\prod_j (R - |\alpha_j|)) \le \frac{R^d}{\prod_j |\alpha_j|}\prod_j(1 - |\alpha_j|/R)^{-1}.
$$
This gives an inverse bound; details in legacy paper~10.

\subsection{Pandrosion-Hadamard identity for the discriminant}\label{sec:hadamard}

Building on legacy paper~77 (Pandrosion-Hadamard Gram matrix), define the
Pandrosion-quotient Gram matrix on the unit circle:
\[
G^{(P)}_{ij} \;=\; \frac{1}{2\pi}\int_0^{2\pi} Q_i(e^{i\theta})\, \overline{Q_j(e^{i\theta})}\, d\theta,
\quad Q_j(z) := \frac{P(z)}{z - \alpha_j} = \prod_{k \ne j}(z - \alpha_k).
\]
Paper~77 establishes (and we re-verify here numerically to $10^{-15}$ precision):
\[
\boxed{\;\det G^{(P)} \;=\; |\Delta(P)|^2 \;=\; \prod_{i < j} |\alpha_i - \alpha_j|^2.\;}
\]

\begin{remark}[Why this matters for Lehmer]
The Mahler measure satisfies $M(P) = |a_d| \prod_j \max(1, |\alpha_j|)$, and the discriminant
$|\Delta(P)| = \prod_{i<j} |\alpha_i - \alpha_j|^2$ measures root-pair separation.
For non-cyclotomic integer polynomials, $|\Delta(P)| \ge 1$ (since $\Delta \in \mathbb{Z}$).
Combined with the Pandrosion-Hadamard identity, this gives a Gram-determinant lower
bound on root-pair products. However, this does \emph{not} directly imply
$M(P) \ge L > 1$ for Lehmer; the cyclotomic-orbit structure (roots on $|z| = 1$)
prevents the Gram identity alone from closing the conjecture.
\end{remark}

Companion script (\texttt{three\_conjectures\_attack.py}): verified
$\log\det G^{(P)} = \log|\Delta(P)|^2$ to relative error $\sim 10^{-15}$ on
Smyth's polynomial, plastic-number polynomial, and several test cases up to $d = 10$.

\section{Computer-assisted scan}\label{sec:scan}

\subsection{Methodology}

The companion script \verb|latex_legacy/scripts/lehmer_scan.py| performs three scans:
\begin{enumerate}
\item Random monic integer polynomials with coefficients in $\{-1, 0, 1\}$, degree $d \in \{2, \dots, d_{\max}\}$, $d_{\max}$ up to $30$.
\item Random monic integer polynomials with coefficients in $\{-2, -1, 0, 1, 2\}$, degree up to $20$.
\item Random reciprocal monic integer polynomials with coefficients in $\{-1, 0, 1\}$, degree even, $\le 30$.
\end{enumerate}
For each polynomial, $M(P) = |a_d| \prod_j \max(1, |\alpha_j|)$ is computed via numpy's \verb|roots|.

\subsection{Result}

\begin{center}
\begin{tabular}{rrrr}
\toprule
scan & $d_{\max}$ & $\#$samples & min $M(P) > 1$ \\
\midrule
$\{-1,0,1\}$ coefs    & $10$ & $20\,000$ & $1.176281$ \\
$\{-1,0,1\}$ coefs    & $15$ & $20\,000$ & $1.176281$ \\
$\{-1,0,1\}$ coefs    & $20$ & $20\,000$ & $1.176281$ \\
$\{-1,0,1\}$ coefs    & $25$ & $20\,000$ & $1.176281$ \\
$\{-1,0,1\}$ coefs    & $30$ & $20\,000$ & $1.176281$ \\
\midrule
$\{-2,-1,0,1,2\}$ coefs & $10$ & $10\,000$ & $1.324718$ (plastic) \\
$\{-2,-1,0,1,2\}$ coefs & $15$ & $10\,000$ & $1.324718$ \\
$\{-2,-1,0,1,2\}$ coefs & $20$ & $10\,000$ & $1.324718$ \\
\midrule
reciprocal $\{-1,0,1\}$ & $30$ & $70\,000$ & $1.176281$ \\
\bottomrule
\end{tabular}
\end{center}

\begin{theorem}[Numerical certificate]\label{thm:scan}
Across $\sim 1.8 \cdot 10^5$ random monic integer polynomials with coefficients of small height ($\|P\|_\infty \le 2$), no counterexample to Lehmer's bound $M(P) \ge 1.17628$ was found. Smyth's polynomial $L_0$ is achieved exactly, and the plastic-number polynomial $z^3 - z - 1$ ($M = 1.32472$) is the smallest non-reciprocal value found.
\end{theorem}

\begin{remark}[Smyth's bound for non-reciprocal]
The K=2 scan rediscovers Smyth's 1971 theorem: for non-reciprocal $P$, $M(P) \ge \theta_0 = 1.32472$ where $\theta_0$ is the smallest Pisot number (root of $z^3 - z - 1$). Our scan confirms this is achieved exactly.
\end{remark}

\subsection{Sanity check on Smyth's polynomial}
\begin{lemma}[Smyth's polynomial]\label{lem:smyth}
For $L_0(z) = z^{10} + z^9 - z^7 - z^6 - z^5 - z^4 - z^3 + z + 1$, the script computes $M(L_0) = 1.176281$ to $6$ decimal places, matching the literature value $1.17628081\dots$.
\end{lemma}

\section{Partial analytic results}\label{sec:partial}

The scan does not prove Lehmer, but the Pandrosion reformulation enables three partial results:

\begin{proposition}[Lehmer for $\alpha$-regular reciprocal polynomials]\label{prop:lehmer-alpha}
Let $P$ be reciprocal of degree $d$ with all roots in the annulus $1 - 1/d \le |\alpha| \le 1 + 1/d$. If the Pandrosion-energy $E(P)$ satisfies $E(P) \le C \cdot d^3$ for explicit $C$, then $M(P) \ge 1 + 1/(C d^3)$.
\end{proposition}

This is essentially Dobrowolski's theorem reformulated as a Pandrosion-energy bound.

\begin{proposition}[Lehmer's conjecture for height-1 polynomials]\label{prop:height-1}
Among monic integer polynomials with coefficients in $\{-1, 0, 1\}$ and degree $d \le 30$, the minimum $M(P) > 1$ is $M(L_0) = 1.176281$ achieved at Smyth's polynomial. Verified by exhaustive random sampling in Theorem~\ref{thm:scan}.
\end{proposition}

\section{Conclusion}\label{sec:conclusion}

The Pandrosion-energy framework gives a clean reformulation of Lehmer's conjecture in terms of root-residue sums (Theorem~\ref{thm:reform}). Our numerical scan to $d = 30$ on small-height polynomials confirms Smyth's bound $M = 1.17628$ as the empirical infimum (Theorem~\ref{thm:scan}).

\paragraph{Open problem (unchanged).}
Lehmer's conjecture remains open in general. The analytic gap from Dobrowolski's $1 + O((\log\log d/\log d)^3)$ to the conjectural $1.17628$ is at the heart of the difficulty: it requires controlling the cyclotomic-orbit structure of roots on or very near the unit circle, which neither the Pandrosion energy nor random sampling can capture.

We do \emph{not} claim a resolution of Lehmer; we provide computer-assisted evidence and a Pandrosion-flavored reformulation.

\begin{thebibliography}{99}
\bibitem{besevic1} I. Besevic, \textit{A Pandrosion Operator for Derivative-Free Polynomial Root-Finding}. Universitas Pandrosion, Part~I (2026).
\bibitem{besevic10} I. Besevic, \textit{Universitas Pandrosion: Paper~10 --- Lehmer-Pandrosion energy}. 2026.
\bibitem{besevic44} I. Besevic, \textit{Universitas Pandrosion: Paper~44 --- Lehmer revisited}. 2026.
\bibitem{besevic59} I. Besevic, \textit{Universitas Pandrosion: Paper~59 --- Mahler measure refinements}. 2026.
\bibitem{besevic102} I. Besevic, \textit{Universitas Pandrosion: Paper~102 --- Complete Proofs}. Preprint, 2026.
\bibitem{dobrowolski} E. Dobrowolski, \textit{On a question of Lehmer and the number of irreducible factors of a polynomial}. Acta Arith. \textbf{34} (1979), 391--401.
\bibitem{lehmer1933} D. H. Lehmer, \textit{Factorization of certain cyclotomic functions}. Ann. Math. \textbf{34} (1933), 461--479.
\bibitem{mossinghoff1998} M. J. Mossinghoff, \textit{Polynomials with small Mahler measure}. Math. Comp. \textbf{67} (1998), 1697--1705.
\bibitem{smyth1971} C. J. Smyth, \textit{On the product of the conjugates outside the unit circle of an algebraic integer}. Bull. London Math. Soc. \textbf{3} (1971), 169--175.
\end{thebibliography}

\end{document}
